VisionServe is functional across all six task types with 16 registered models, but
several known limitations remain.

\begin{itemize}

    \item \textbf{Pending model exports.}
          The RT-DETR and Depth-Anything-V2 ONNX exports hosted by
          \texttt{onnx-community} on Hugging Face currently return HTTP 401 and require
          gated authentication; end-to-end benchmarks for these models are therefore not
          yet available. NanoSAM~\cite{zhang2023mobilesam} requires a manual ONNX export
          from the NVIDIA-AI-IOT/nanosam repository via Google Drive, which cannot be
          automated in the existing manifest-based download pipeline.

    \item \textbf{CLIP text encoder deferred.}
          Only the CLIP~\cite{radford2021clip} image encoder is currently implemented.
          Integrating the text encoder requires a byte-pair encoding (BPE) tokenizer,
          which has been deferred to v2. As a result, zero-shot classification is
          currently limited to scenarios where text embeddings are precomputed externally
          and supplied as fixed vectors.

    \item \textbf{No ROCm execution provider.}
          The \texttt{yalue/onnxruntime\_go} binding does not yet expose the ROCm
          execution provider. AMD GPU users must therefore fall back to CPU inference,
          foregoing GPU acceleration on that platform.

    \item \textbf{Simplified automatic mask generation.}
          MobileSAM~\cite{zhang2023mobilesam} supports prompt-free automatic mask
          generation (AMG), but the implementation is a single-resolution approximation
          of SAM's full pipeline: it tiles the image with a fixed $16\times16$ grid of
          foreground point prompts (256 decoder calls sharing one encoder pass), keeps
          masks whose predicted IoU exceeds a fixed threshold ($0.85$) and whose area
          falls within heuristic bounds ($10^{-4}$ to $0.95$ of the image), then
          deduplicates by greedy pixel-IoU NMS ($0.70$). It does \emph{not} implement
          SAM's multi-crop / multi-scale AMG, stability-score filtering, or box-NMS
          across crops, so dense small-object recall is lower than the reference
          generator. EfficientSAM~\cite{xiong2023efficientsam} and SAM2~\cite{ravi2024sam2}
          remain prompt-only.

    \item \textbf{Single-image batching only.}
          The server currently processes one image per request (\texttt{batch\_size=1}).
          Submitting multiple images in a single call is not supported, which limits
          throughput in offline or batch-processing scenarios where GPU utilisation could
          otherwise be improved through request batching~\cite{nvidia2023triton}.

    \item \textbf{Windows support untested.}
          The Go toolchain produces a valid Windows binary, and the DirectML execution
          provider is listed in the EP allowlist, but end-to-end operation on Windows
          has not been validated. Windows-specific path handling and ORT library loading
          may require additional work.

    \item \textbf{Low-load latency regression.}
          At concurrency~1 VisionServe is slower per request than the GIL-bound
          FastAPI+ONNX~Runtime baseline and far slower than Triton: for
          MobileNetV3-Small its p50/min latency is $57.2$/$25.9$\,ms versus FastAPI's
          $16.1$/$10.9$\,ms and Triton's $4.0$/$2.9$\,ms. The cost is not the GPU
          (ORT~Run~$<$1\,ms) but pure-Go JPEG decode ($\sim$19.2\,ms): Pillow uses
          SIMD libjpeg-turbo and Triton receives a pre-decoded tensor. A build-tagged
          cgo libjpeg-turbo would close the gap but breaks the single-static-binary
          property, so we keep pure-Go decode as a deliberate trade-off and report the
          $\Delta$.

    \item \textbf{Concurrency advantage scoped to GIL-bound CPython.}
          The throughput crossover (C2) is measured against the deployed-default
          CPython interpreter with the global interpreter lock. Free-threaded (no-GIL)
          Python~3.13 narrows this gap, so the crossover is scoped to GIL-bound CPython
          and \emph{not} claimed against free-threaded Python.

    \item \textbf{License enforcement trusts the declared license.}
          The load-time guarantee (C1) is license-\emph{policy} enforcement, hardened by
          content hashing (sha256) and a verified-source allowlist; it does \emph{not}
          semantically verify that a model's declared license matches its true legal
          status. By binding declared license $\leftrightarrow$ weight bytes
          $\leftrightarrow$ audited source, it mitigates relicensing and
          permissive-washing, but a deliberately mislabeled upstream license is out of
          scope of the threat model.

    \item \textbf{\texttt{duration\_ms} not directly comparable to the baseline.}
          VisionServe's \texttt{duration\_ms} covers preprocess~+~ORT~+~postprocess,
          whereas the FastAPI baseline's server timer covers ORT only. Raw
          \texttt{duration\_ms} is therefore not like-for-like; end-to-end client latency
          is the comparable metric.

    \item \textbf{Triton comparison is default-config.}
          The Triton parity result is against an out-of-the-box Triton (HTTP, single
          model instance, no dynamic batching). A tuned Triton (gRPC + shared memory +
          dynamic batching) would exceed it; no parity is claimed against a tuned
          Triton.

\end{itemize}
